\section{Conclusion and Research Programme}
\label{sec:conclusion}

\subsection{What this paper establishes}
\label{sec:conclusion-established}

This paper has proposed a framework that predicts, from a language's grammar alone, how much of a transformer's parameter budget morphology-aware tokenization saves over a byte-pair baseline. The two ingredients are the Shannon entropy $H(L)$ of feature bundles a word can carry, and the structural recoverability coefficient $\rho(L)$ that captures how much of that entropy a surface-only parser can recover. Both are read off a morphological grammar engine before any training run begins. Equation~\ref{eq:rebate} turns them into a quantitative prediction: how much grammatical structure the model no longer has to learn distributionally because the tokeniser already exposes it.

Tested on the Phase~1 six-model experiment that paired BPE and morphology-aligned variants for English, Arabic, and Turkish (analytic, templatic, and agglutinative in turn), the framework's ranking prediction held. The observed rebate is monotonically ordered by $\rho(L) \cdot H(L)$, a structural quantity derivable from the grammar alone. The magnitude prediction did not. Under any reasonable parameterisation of the Hoffmann scaling-law derivative, the implied per-bit coefficient $\alpha$ varies by roughly forty-fold across the three languages. We named this empirical pattern the \textbf{Agglutinative Compounding Effect} and argued that two mutually exclusive readings are consistent with Phase~1 evidence: either the true rebate is genuinely super-linear in $\rho \cdot H$ (Interpretation~A), or the standard scaling-law machinery applies only approximately in the Chinchilla-suboptimal regime Phase~1 occupied (Interpretation~B).

The three contributions the paper claims are therefore the framework itself, the ranking result it predicted correctly, and the named phenomenon the framework did not anticipate.

\subsection{What has been built since the original draft}
\label{sec:conclusion-built}

Several items the original draft listed as follow-ups have since been built and integrated into the architecture documented in \S\ref{sec:framework-engine}. These changes were carried out before the experimental phase scoped in \S\ref{sec:conclusion-programme} below, and they shape what the experimental phase can now ask.

\paragraph{The 3rd-tier wazn-semantic-class layer for Arabic, built.}
The wazn-semantic-class extension named below as a follow-up has been implemented. The Arabic engine's Step B' propagates the \texttt{semantic\_role} field of every matched template (185 distinct values across 348 awz\=an in \texttt{ar\_templates.json}) into the token's tag bundle and emits it as its own input stream alongside the existing surface, feature-bundle, and root streams. This is the 4-stream Arabic architecture the framework now actually trains: $x = \mathrm{tok\_emb} + \mathrm{feat\_emb} + \mathrm{root\_emb} + \mathrm{wazn\_emb}$ rather than the 3-stream architecture the original 2nd-tier experiment used. The experimental phase below therefore tests not the 2nd-tier vs.\ 1st-tier comparison the original draft scoped, but the 4-stream vs.\ baseline comparison the post-everything engine permits.

\paragraph{A sentence-grammar layer for all four languages.}
A layer the original draft did not envisage was added between per-word morphology and downstream consumption: a sentence-grammar package per language that splits text on punctuation, walks each sentence resolving POS ambiguities by neighbour context, and validates the sentence against a small set of allowed structural templates encoded as Python predicates with rejection messages in each language's own grammatical terminology. The morphology proposes; the grammar disposes. Architecturally, this means the framework's input signal is no longer a per-word function of the surface but a sentence-window function, which closes a gap reviewers might otherwise have asked about. The same layer underpins the Arabic nahwiyya guards (nawasikh, idafa, hal-clause agreement) and the Mandarin BA/BEI-construction guards that would have been left implicit in the original draft.

\paragraph{Mandarin as a fourth language.}
The original Phase~1 covered English, Arabic, and Turkish. The post-everything project covers four languages by adding Mandarin as a structural outlier (isolating, with the Kangxi radical-class as the analog of Arabic's wazn-semantic axis at a smaller scale). The Mandarin engine, its 214-class radical lookup, and its 2-stream architecture (character + radical-class) are documented in \S\ref{sec:cases-zh}. The experimental phase below therefore runs four head-to-head comparisons, not three.

\subsection{Research programme}
\label{sec:conclusion-programme}

Two experiments would discriminate between Interpretation~A and Interpretation~B, and between either of them and the possibility that the framework is wrong in some third way not yet articulated. Both experiments are within reach of the project's existing infrastructure.

\paragraph{Spectrum expansion at the current scale.}
Five additional grammar engines, for German (fusional), Spanish (regular fusional), Hungarian (agglutinative-fusional), Swahili (classificatory), and Basque (polypersonal, ergative, isolate), were built and tested in earlier rounds of the project and remain in archived state. Running the same pipeline across these five engines would produce nine data points spanning the full morphological spectrum plus three structural outliers. Under Interpretation~A, the super-linear shape observed in the four-language case should extend smoothly. Under Interpretation~B, the language-dependent distortion from Chinchilla-suboptimal training should persist across the expanded set but without a typologically-meaningful pattern. Under a framework-failure reading, the ordinal match of \S\ref{sec:results-ordinal} should degrade.

\paragraph{Scale ladder within the existing four languages.}
Training each of the eight current models (4 languages $\times$ 2 regimes) at small ($20 \times 10^6$ parameters, $10^8$ tokens) and medium ($6 \times 10^7$ parameters, $1.2 \times 10^9$ tokens) scales would place the loss-to-parameter derivative in a well-characterised Chinchilla-optimal regime. Under Interpretation~A, the super-linearity should persist across scales. Under Interpretation~B, it should vanish or substantially compress as training approaches optimality. The infrastructure for this ladder is in place: the multi-stream MiniGPT architecture, the tokenisation pipeline, and the SageMaker launcher are documented in the project repository and verified end-to-end on smoke-scale data.

Together the two experiments produce a $2 \times 2$ evidence grid against which both interpretations are falsifiable. Neither experiment is conceptually novel; both are bounded in compute cost; both are specified by the framework at the level of concrete engines to be run and metrics to be compared.

\paragraph{Low-resource edge-deployment showcase for templatic languages.}
Arabic's root-based lexicon is vocabulary-efficient: roughly 10{,}000 triconsonantal roots generate more than 100{,}000 distinct surface word forms. A tokeniser that splits words into root and pattern can cover Arabic with a vocabulary an order of magnitude smaller than a conventional BPE tokeniser of equivalent coverage. A targeted extension would train a small ($\sim 2 \times 10^6$ parameter) Arabic model with a radically capped vocabulary under the 1st-tier morph regime of \S\ref{sec:results}, and compare its downstream performance (part-of-speech tagging, dialect identification, natural-language inference) against a BPE baseline of matched parameter count and matched vocabulary budget. A win for the morph model at small vocabulary budget would directly instantiate the framework's claim that morphology-aware tokenisation makes templatic languages newly hostable on devices the size of a smartphone, and would carry a concrete equity argument: the languages most under-served by current pipelines are exactly the ones most rewarded by this approach.

\paragraph{Per-model loss attribution on the Arabic compositional bucket.}
\S\ref{sec:results-comp-gen} constructed a 4{,}747-instance HELD\_OUT bucket on the Arabic test corpus made of word forms whose (root, template) combination is unseen in training, even though both the root and the template are individually attested. A clean per-model evaluation requires assigning per-token loss back to its source word and then aggregating only over positions whose source word lies in each bucket. This in turn requires keeping word-to-token alignments separately for the BPE tokeniser, the shallow morph tokeniser, and the engine-root tokeniser. With that bookkeeping in place, the test would compare $\Delta\mathcal{L}$(morph $-$ baseline) on HELD\_OUT against the same on SEEN\_PAIR. A larger morph advantage on HELD\_OUT would be direct evidence that morphology-aware tokenisation induces compositional learning of Arabic's root--template system, as opposed to memorisation of surface-and-bundle co-occurrence.

\paragraph{Cross-Semitic transfer of the Arabic engine structure.}
Arabic's root-and-pattern morphology is a property the whole Semitic family shares: Hebrew, Aramaic, Amharic, Maltese, and others. The Arabic grammar engine of \S\ref{sec:cases-ar} is in principle portable to related languages with modest changes to the root extractor and the template inventory. A lighter-weight follow-up would port the engine to Hebrew or Maltese and re-compute $H(L)$, $\rho(L)$, and $\rho \cdot H$ without training any models, showing that the framework's pre-training predictive signal extends across the Semitic family with negligible additional infrastructure. This reach matters because the framework's practical case is strongest exactly where the digital infrastructure is weakest, and the Semitic family contains several languages in that position.

\subsection{Stance}
\label{sec:conclusion-stance}

The paper's position is straightforward. A theoretical framework that ranks three typologically distant languages correctly while failing to predict their magnitudes under any current scaling law is more scientifically productive than either a confirmatory gradient result or a scaling-law extension with no structural grounding. The Agglutinative Compounding Effect, whether it turns out to be a real property of morphology-aware language modelling or dissolves into a correction term on existing scaling laws, is the question the framework was built to ask. The research programme of \S\ref{sec:conclusion-programme} is what asking that question completely would look like.
